%document
\documentclass[10pt]{beamer}
%theme
\usetheme{CambridgeUS}
% packages
\usepackage{color}
\usepackage{listings}
\usepackage[ngerman]{babel}
\usepackage[utf8]{inputenc}
\usepackage{multicol}


% color definitions
\definecolor{mygreen}{rgb}{0,0.6,0}
\definecolor{mygray}{rgb}{0.5,0.5,0.5}
\definecolor{mymauve}{rgb}{0.58,0,0.82}

\lstset{
    backgroundcolor=\color{white},
    % choose the background color;
    % you must add \usepackage{color} or \usepackage{xcolor}
    basicstyle=\footnotesize\ttfamily,
    % the size of the fonts that are used for the code
    breakatwhitespace=false,
    % sets if automatic breaks should only happen at whitespace
    breaklines=true,                 % sets automatic line breaking
    captionpos=b,                    % sets the caption-position to bottom
    commentstyle=\color{mygreen},    % comment style
    % deletekeywords={...},
    % if you want to delete keywords from the given language
    extendedchars=true,
    % lets you use non-ASCII characters;
    % for 8-bits encodings only, does not work with UTF-8
    frame=single,                    % adds a frame around the code
    keepspaces=true,
    % keeps spaces in text,
    % useful for keeping indentation of code
    % (possibly needs columns=flexible)
    keywordstyle=\color{blue},       % keyword style
    % morekeywords={*,...},
    % if you want to add more keywords to the set
    numbers=left,
    % where to put the line-numbers; possible values are (none, left, right)
    numbersep=5pt,
    % how far the line-numbers are from the code
    numberstyle=\tiny\color{mygray},
    % the style that is used for the line-numbers
    rulecolor=\color{black},
    % if not set, the frame-color may be changed on line-breaks
    % within not-black text (e.g. comments (green here))
    stepnumber=1,
    % the step between two line-numbers.
    % If it's 1, each line will be numbered
    stringstyle=\color{mymauve},     % string literal style
    tabsize=4,                       % sets default tabsize to 4 spaces
    % show the filename of files included with \lstinputlisting;
    % also try caption instead of title
    language = Java,
	showspaces = false,
	showtabs = false,
	showstringspaces = false,
	escapechar = ,
}

\def\ContinueLineNumber{\lstset{firstnumber=last}}
\def\StartLineAt#1{\lstset{firstnumber=#1}}
\let\numberLineAt\StartLineAt



\newcommand{\codeline}[1]{
	\alert{\texttt{#1}}
}

% This Document contains the information about this course.

% Authors of the slides
\author{Vincent Gerber, Tilman Hinnerichs}

% Name of the Course
\institute{Java Kurs}

% Fancy Logo
\titlegraphic{\hfill\includegraphics[height=1.25cm]{../templates/fsr_logo_cropped}}


\title{Java}
\subtitle{Introduction}
\date{\today}


\begin{document}

% \section{Organisation}
\begin{frame}
	\titlepage
\end{frame}
\begin{frame}{Overview}
	\setbeamertemplate{section in toc}[sections numbered]
	\tableofcontents
\end{frame}

\section{Proceeding}
\begin{frame}{About this course}
% 	Language?\\
	Requirements
	\begin{itemize}
		\item You know how to use a computer
        \item If possible, bring your own computer with you
		\item Probably knowledge of other programming languages
% 		\item Java is not your first programming language or you are a fast learner
	\end{itemize}
	Proceeding
	\begin{itemize}
		\item There will be ~14 lessons
		\item Each covers a topic and comes with excercises
	\end{itemize}
\end{frame}

% \subsection{Resources}
\begin{frame}{Some resources}
	\begin{itemize}
		\item You can ask your tutor
		\item Join the Auditorium group \hfill \\
			\url{http://auditorium.inf.tu-dresden.de}
		\item StackOverflow, FAQs, Online-tutorials, ... \hfill \\
		\item Official documentation \hfill \\
			\url{https://docs.oracle.com/javase/8/}
        \item mailinglist \url{programmierung@ifsr.de}
        %\item Cyberspace (wednesday 5./6. DS)
		\item Material-Repository \\
			\url{https://github.com/THinnerichs/Java-Course}
	\end{itemize}
\end{frame}

\begin{frame}{Why is there a Java course?}
	\begin{itemize}
		\item Help for your SWT lecture
		\item Help for your SWT project next term
		\item Java is easy to start with, if you haven't worked with OOP before
		\item Tons of jobs for Java developers
		\item[] \textbf{AND} for your questions and ideas for the course
	\end{itemize}
\end{frame}

\begin{frame}[c]{}
	\centering
	\textsc{\Huge{\textbf{Questions?}}}
\end{frame}

\begin{frame}{About Java}
	% \tikzoverlay at (7cm,1.4cm) {
	% 	\includegraphics[width=3cm, height=3cm]{res/logo-java.png}
	% }
	Pros:
	\begin{itemize}
		\item Syntax like C++
		\item Strongly encourages OOP
		\item Platform-independent (JVM)
		\item Very few external libraries
		\item[] $\Rightarrow$ Easy to use and very little to worry about
	\end{itemize}
\end{frame}

\begin{frame}{About Java}
	% \tikzoverlay at (7cm,1.4cm) {
	% 	\includegraphics[width=3cm, height=3cm]{res/logo-java.png}
	% }
	Cons:
	\begin{itemize}
		\item A lot of unnecessary features \\ in the JDK
		\item Slower than assembly
		\item No multi-inheritance
		\item Weak generics
		\item Mediocre support for other programming paradigms
		\item[] $->$ Neither fast, small nor geeky
	\end{itemize}
\end{frame}

\section{Your first program}
\begin{frame}{Hello World}
  DEMO
\end{frame}

\begin{frame}[fragile]{Creating your Working Environment}
  Open the Terminal
  \begin{lstlisting}
      mkdir myProgram
      cd myProgram
      touch Hello.java
      vim Hello.java
  \end{lstlisting}
\end{frame}

\subsection{Hello World!}

\begin{frame}[fragile]{Hello World!}
	This is an empty JavaClass.
  Java Classes always start with a capital letter
	\begin{lstlisting}
	public class Hello {

	}
	\end{lstlisting}
\end{frame}

\begin{frame}[fragile]{Hello World!}
	This is a small program printing \emph{Hello World!} to the console:
	\begin{lstlisting}
	public class Hello {
	    public static void main(String[] args) {
	        System.out.println("Hello World!");
	    }
	}
	\end{lstlisting}
\end{frame}

\begin{frame}[fragile]{How to run your program}
   save your program by pressing 'esc', then ':w'
   exit vim by typing ':q' (and hit return)
   then:
   \begin{lstlisting}[language=bash]
      javac Hello.java
      java Hello
   \end{lstlisting}
\end{frame}

\subsection{Setting up IntelliJ IDEA}

\begin{frame}{Hello World in an IDE}
   DEMO
\end{frame}

\begin{frame}{Receive a copy of IntelliJ IDEA}
  IntelliJ IDEA is a powerful IDE\footnote{Integrated Development Environment}, e.g. for Java.
  \begin{itemize}
      \item You can download IntelliJ IDEA at  \\
          \url{https://www.jetbrains.com/idea/}
      \item Get an Ultimate-License at \\
          \url{https://www.jetbrains.com/student/}
      \item Use JetBrains IDEs for all programming languages
  \end{itemize}
  Ecplipse is free and open-source, but less powerful.
\end{frame}

\section{Basics}

\begin{frame}[fragile]{Comments}
   \begin{lstlisting}
   public class Hello {
       // prints a "Hello World!" on your console
       public static void main(String[] args) {
           System.out.println("Hello World!");
       }
   }
   \end{lstlisting}
   You should always comment your code. \\
   Code is read more often than it is written.
   \begin{itemize}
       \item // single line comment
       \item /* comment spanning \\
           multiple lines */
   \end{itemize}
\end{frame}

\subsection{Some definitions}

\begin{frame}[fragile]{Code concepts}
	\begin{lstlisting}
	public class Hello {
	    // Calculates some stuff and outputs everything on the console
	    public static void main(String[] args) {	        
	        int x;
	        x = 9;
	        int y = 23;
	        int z;
	        z = x * y;
	        
	        System.out.println(z);
	    }
	}
	\end{lstlisting}
\end{frame}

\begin{frame}[fragile]{Code concepts}
	\begin{lstlisting}
	public class Hello {
	    // Calculates some stuff and outputs everything on the console
	    public static void main(String[] args) {
	        System.out.println(9 * 23);
	    }
	}
	\end{lstlisting}
\end{frame}

\begin{frame}{Primitive data types}
	Java supports some primitive data types:
	\begin{itemize}
		\item[boolean] a truth value (either \textbf{true} or \textbf{false})
		\item[int] a 32 bit integer
		\item[long] a 64 bit integer
		\item[float] a 32 bit floating point number
		\item[double] a 64 bit floating point number
		\item[char] an ascii character
		\item[void] the empty type (needed in later topics)
	\end{itemize}
\end{frame}

\begin{frame}[fragile]{About the Semicolon}
	\begin{lstlisting}
	public class Hello {
	    // prints a "Hello World!" on your console
	    public static void main(String[] args) {
	        System.out.println("Hello World!")@\textcolor{red}{\texttt{;}}@
	    }
	}
	\end{lstlisting}
	Semicolons conclude all statements. \\
	Blocks do not need a semicolon.
\end{frame}

\begin{frame}[fragile]{Blocks}
	\begin{lstlisting}
	public class Hello @\textcolor{red}{\texttt{\{}}@
	    // prints a "Hello World!" on your console
	    public static void main(String[] args) {
	        System.out.println("Hello World!");
	    }
	@\textcolor{red}{\texttt{\}}}@
	\end{lstlisting}
	Everything between \{ and \} is a \emph{block}. \\
	Blocks may be nested.
\end{frame}

\begin{frame}[fragile]{Naming of Variables}
	\begin{itemize}
		\item The names of variables can begin with any letter or underscore. \\
		Usually the name starts with small letter.
		\item Compound names should use CamelCase.
		\item Use meaningful names.
	\end{itemize}
	\begin{lstlisting}
	public class Calc {
	    public static void main(String[] args) {
	    	int a = 0; // not very meaningful
	    	float myFloat = 5.3f; // also not meaningfull
	    	int count = 7; // quite a good name

	    	int rotationCount = 7; // there you go
	    }
	}
	\end{lstlisting}
\end{frame}

\subsection{Calculating}

\begin{frame}[fragile, allowframebreaks]{Calculating with \emph{int}}
	\begin{lstlisting}
	public class Calc {
	    public static void main(String[] args) {
	        int a; // declare variable a
	        a = 7; // assign 7 to variable a
	        System.out.println(a); // prints: 7
	        a = 8;
	        System.out.println(a); // prints: 8
	        a = a + 2;
	        System.out.println(a); // prints: 10
	    }
	}
	\end{lstlisting}
	After the first assignment the variable is initialized.
\framebreak
	\begin{lstlisting}
	public class Calc {
	    public static void main(String[] args) {
	        int a = -9; // declaration and assignment of a
	        int b; // declaration of b
	        b = a; // assignment of b
	        System.out.println(a); // prints: -9
	        System.out.println(b); // prints: -9
	        a++; // increments a
	        System.out.println(a); // prints: -8
	    }
	}
	\end{lstlisting}
% \framebreak
% 	\begin{lstlisting}
% 	public class Calc {
% 	    public static void main(String[] args) {
% 	        int b; // declaration of b
% 	        System.out.println(b);
% 	    }
% 	}
% 	\end{lstlisting}
% 	Uninitialized variables will cause an Exception. \\
% 	An Exception is a kind of error we will discuss later.\\
% 	\vspace{1em}
% 	\emph{Always assign your variables!}
% \framebreak

\framebreak
	Some basic mathematical operations:
	\begin{tabular}{ll}
		Addition & \texttt{a + b;} \\
		Subtraction & \texttt{a - b;} \\
		Multiplication &\texttt{a * b;} \\
		Division & \texttt{a / b;} \\
		Modulo & \texttt{a \% b;} \\
		Increment & \texttt{a++;} \\
		Decrement & \texttt{a--;} \\
	\end{tabular}
\end{frame}

\begin{frame}[fragile, allowframebreaks]{Calculating with \emph{float}}
	\begin{lstlisting}
	public class Calc {
	    public static void main(String[] args) {
	        float a = 9;
	        float b = 7.5f;
	        System.out.println(a); // prints: 9.0
	        System.out.println(b); // prints: 7.5
	        System.out.println(a + b); // prints: 16.5
	    }
	}
	\end{lstlisting}
\framebreak
	\begin{lstlisting}
	public class Calc {
	    public static void main(String[] args) {
	        float a =       8.9f;
	        float b = 3054062.5f;
	        System.out.println(a); // prints: 8.9
	        System.out.println(b); // prints: 3054062.5
	        System.out.println(a + b); // prints: 3054071.5
	    }
	}
	\end{lstlisting}
	Float has a limited precision. \\
	\emph{This might lead to unexpected results!}
\end{frame}

\begin{frame}[fragile]{Mixing \emph{int} and \emph{float}}
	\begin{lstlisting}
	public class Calc {
	    public static void main(String[] args) {
	        float a = 9.3f;
	        int b = 3;
	        System.out.println(a + b); // prints: 12.3
	        float c = a + b;
	        System.out.println(c); // prints: 12.3
	    }
	}
	\end{lstlisting}
	Java converts from \textbf{int} to \textbf{float} by default, if necessary. \\
	But not vice versa.
\end{frame}

\subsection{Text with Strings}

\begin{frame}[fragile]{Strings}
	A String is not a primitive data type but an object. \\
	We discuss objects in detail in the next section.
	\begin{lstlisting}
	public class Calc {
	    public static void main(String[] args) {
	        String hello = "Hello World!";
	        System.out.println(hello); // print: Hello World!
	    }
	}
	\end{lstlisting}
\end{frame}

\begin{frame}[fragile]{Concatenation}
	\begin{lstlisting}
	public class Calc {
	    public static void main(String[] args) {
	        String hello = "Hello";
	        String world = " World!";
	        String sentence = hello + world;
	        System.out.println(sentence);
	        System.out.println(hello + " World!");
	    }
	}
	\end{lstlisting}
	You can concatenate Strings using the +. Both printed lines look the same.
\end{frame}

\begin{frame}[fragile]{Strings and Numbers}
	\begin{lstlisting}
	public class Calc {
	    public static void main(String[] args) {
	    	int factorA = 3;
	    	int factorB = 7;
	    	int product = factorA * factorB;
	    	String answer =
	            factorA + " * " + factorB + " = " + product;
	        System.out.println(answer); // prints: 3 * 7 = 21
	    }
	}
	\end{lstlisting}
	Upon concatenation, primitive types will be replaced by their current value as \emph{String}.
\end{frame}

\begin{frame}[fragile]{Summary}
	What we learned today: \\
	\begin{itemize}
		\item How to write our first program
		\item How to execute our first program
		\item How to do basic calculations
		\item How to work with strings
	\end{itemize}
\end{frame}

\end{document}
